\chapter{Introduction}
\label{chap:introduction}

\section{Background \& Problem}

Low birth weight (LBW), defined as a birth weight less than 2.5 kilograms \citep{lbw_def, kramer1987}, is a significant public health indicator. Infants born with LBW face substantially higher risk for neonatal mortality, developmental delays, and chronic health problems such as respiratory and neurological impairments \citep{finch2003}. These adverse outcomes arise from a complex interplay of genetic, biological, environmental, and socioeconomic factors. Understanding and identifying determinants of LBW is therefore critical for developing targeted preventative measures and improving neonatal outcomes. 

An early landmark meta-analysis by \citet{kramer1987} reviewed 895 studies from 1974-1984 and identified 43 causal determinants of LBW. Kramer concluded that maternal anthropometry (height and pre-pregnancy weight), inadequate gestational weight gain, cigarette smoking, malaria infection, and a history of adverse pregnancy outcomes exert \emph{independent} effects on intrauterine growth restriction (IUGR), while few factors influence gestational duration \citep{kramer1987}. The highly-interrelated nature of these risk factors leads to confounding and presents a significant challenge in modeling birth-weight outcomes by their interactive, not additive, effects. For example, inadequate weight gain during pregnancy might depend on whether a mother smokes or has underlying health conditions. Even carefully specified regression models, match machine learning performance, but often struggle to disentangle which combinations of maternal-infant characteristics \emph{truly} signify an at-risk pregnancy. This motivates the use of nonparametric, tree-based methods that enable the identification of complex interactions among risk factors, allowing for more nuanced modeling of birth-weight outcomes.

% Decades of research confirm that LBW is a multifactorial issue. An early landmark meta-analysis by \citet{kramer1987} reviewed 895 studies from 1974-1984 and identified 43 causal determinants. Kramer concluded that maternal anthropometry (height and pre-pregnancy weight), inadequate gestational weight gain, cigarette smoking, malaria infection, and a history of adverse pregnancy outcomes exert \emph{independent} effects on intrauterine growth restriction (IUGR), while few factors influence gestational duration \citep{kramer1987}. The highly-interrelated nature of these risk factors led to confounding, yielding an impediment for modeling birth-weight outcomes by their interactive, not additive, effects. For example, inadequate pregnancy weight gain might depend on whether she smokes or has health conditions. Classical regression models such as logistic regression, typically assume additive effects and thus miss such interactions. As a result, traditional models often struggle to disentangle which combinations of maternal-infant characteristics \emph{truly} signify an at-risk pregnancy. 

% (STENGTHEN THIS PARAGRAPH: "The highly-interrelated nature of these risk factors led to confounding, yielding an impediment for modeling birth-weight outcomes by their interactive, not additive, effects. For example, inadequate pregnancy weight gain might depend on whether she smokes or has health conditions. Classical regression models such as logistic regression, typically assume additive effects and thus miss such interactions. As a result, traditional models often struggle to disentangle which combinations of maternal-infant characteristics \emph{truly} signify an at-risk pregnancy." with a sentence about how this motivates the use of tree-based methods.)

Subsequent work in epidemiology repeatedly show that these risk factors are not found in isolation. In \citeyear{KITSANTAS2006275}, \citet{KITSANTAS2006275} used Classification and Regression Trees (CART), developed by \citet{breiman1984classification}, to identify high-risk profiles for LBW a large birthrecord dataset. They found that mothers who smoked \emph{and} had inadequate weight-gain during pregnancy had sharply elevated LBW risk. This study highlights the strong interpretability of CART deployed to a high-dimensional dataset, but still lacks predictive power over logistic regression by relying entirely on empirical observations. Moreover, reducing the problem to a binary LBW indicator variable discards vital information of the severity of LBW outcomes. Despite differences, a baby born far below 2.5 kg is treated the same as a baby just below the threshold. This binary cutoff masks important differences in the birth-weight distribution, and limits the ability to identify risk factors that influence the severity of LBW outcomes.

% They found that mothers who smoked \emph{and} had inadequate weight-gain during pregnancy had sharply elevated LBW risk. This study highlights the strengths of interpretability using CART, but still lacks predictive power over logistic regression by relying entirely on empirical observations. Moreover, reducing the problem to a binary LBW indicator variable discards vital information about how far a newborn falls below the LBW threshold. Regardless of differences, a baby just above 2.5 kg is treated the same as a much heavier baby, and all LBW cases are treated alike. This binary cutoff masks important differences in the birth-weight distribution.

Recent research has moved beyond classification toward estimating the full birth-weight distribution conditional on covariates. Bayesian nonparametric mixture models allow the birth-weight density to vary flexibly across subpopulations defined by maternal factors, without strong parametric assumptions \citep{dunson2008}. Other approaches use copula-based or density regression techniques to jointly model birth weight with related outcomes such as gestational age \citep{rathjens2023}. (ADD ANY OTHERS) These methods can capture detailed distributional effects of predictors, such as how covariates influence the entire left tail of the birth-weight distribution. However, a drawback in these advanced models is their complexity and lack of interpretability for users. In contrast, practitioners particularly in public health often prefer models that yield clear, simple decision rules for identifying high-risk subgroups. 

Several demographic and socioeconomic factors are well-known to influence LBW risk. Mothers who give birth at younger or older ages, have lower educational attainment, or experience limited access to prenatal care are at elevated risk for LBW outcomes \citep{age_differences_lbw, finch2003, jain2024}. Environmental exposures, including tobacco use, substance abuse, and air pollution, further elevate risk by interfering with fetal growth and development \citep{standford_med_lbw, lu2020combined}. Crucially, incidence of LBW newborns vary radically by racial and economic demographics. In the US, this incidence rate for Black infants is double that of caucasian newborns, at 14.7\% to 7.1\%  \citep{marchofdimes2024}. These patterns underscore the multifactorial nature of effects that influence LBW outcomes and the need to account for diverse influences in any predictive model.

% Younger and older mothers are associated with higher incidence of LBW \citep{age_differences_lbw}. Lower educational attainment is associated with limited health care access \citep{finch2003,jain2024}, and environmental exposures, including tobacco use, substance abuse, and air pollution, further elevate risk by interfering with fetal growth and development \citep{standford_med_lbw, lu2020combined}. Inadequate prenatal care is another important factor of LBW outcomes \citep{prenatal_lbw}. Crucially, LBW incidence also varies sharply by racial and economic contexts. In the United States, the LBW incidence rate for Black infants is double that of white newborns, comparing 14.7\% to 7.1\%  \citep{marchofdimes2024}. These patterns underscore the multifactorial nature of effects that influence LBW outcomes and the need to account for diverse influences in any predictive model. 

Taken together, these considerations highlight a central challenge in birth-weight modeling: existing methods trade flexibility for interpretability. Approaches using decision trees alone provide transparent subgroup rules while ignoring the full birth-weight distribution, simultaneously advanced Bayesian density models capture distributional details but lack intuitive clarity. This work aims to bridge the gap by developing a quasi-Bayesian tree-based framework that stratifies the population into interpretable risk subpopulations while modeling full birth-weight distributions for predicting LBW outcomes.


% FIGURES -------------------------------------------------------
% Figures
% ranked improvement - FULL
\begin{figure}[H]
    \centering
    \includegraphics[width=1\linewidth]{amkurth-Thesis/chapter3/figures/improvement/tree_split_contribution_top20_Full Model.pdf}
    \caption[Full Model Ranked Improvement]{Full Model Ranked Improvement. Rankings represent summed reduction in deviance (improvement in model fit) across all nodes where each variable is used for splitting in the tree. Plot only displays top 20 ranked variables.}
    \label{fig:full-model-ranked-imp}
\end{figure}

% ranked improvement - LBW
\begin{figure}[H]
    \centering
    \includegraphics[width=1\linewidth]{amkurth-Thesis/chapter3/figures/improvement/tree_split_contribution_top9_LBW-only Model.pdf}
    \caption[LBW-Only Model Ranked Improvement]{LBW-Only Model Ranked Improvement. Rankings represent summed reduction in deviance (improvement in model fit) across all nodes where each variable is used for splitting in the tree. Plot displays all ranked variables.}
    \label{fig:lbw-model-ranked-imp}
\end{figure}

% % First barplot figure
% \begin{figure}[H]
%     \centering
%     \includegraphics[width=1\textwidth]{amkurth-Thesis/chapter3/figures/high_low_risk_full.pdf}
%     \caption{Full model: Aggregated mean probability estimates by birth-weight category across $B$ bootstraps, showing the distribution of predicted probabilities for high and low risk subgroups with 95\% percentile intervals.}
%     \label{fig:high-low-risk-full}
% \end{figure}

% % First barplot figure (without NBW)
% \begin{figure}[H]
%     \centering
%     \includegraphics[width=1\textwidth]{amkurth-Thesis/chapter3/figures/high_low_risk_full_without_NBW.pdf}
%     \caption{Full model (without NBW birth-weight column): Aggregated mean probability estimates by birth-weight category across $B$ bootstraps, showing the distribution of predicted probabilities for high and low risk subgroups with 95\% percentile intervals.}
%     \label{fig:high-low-risk-full-w/o-NBW}
% \end{figure}

% % Second barplot figure
% \begin{figure}[H]
%     \centering
%     \includegraphics[width=1\textwidth]{amkurth-Thesis/chapter3/figures/high_low_risk_small.pdf}
%     \caption{LBW-only model: Aggregated mean probability estimates by birth weight category across $B$ bootstraps, showing the distribution of predicted probabilities for high and low risk subgroups with 95\% percentile intervals.}
%     \label{fig:high-low-risk-lbw}
% \end{figure}

% ─────────────────────────────────────────────────────────────────────────
% Full + NBW (11 categories)
\begin{figure}[H]
  \centering
  \includegraphics[width=\textwidth]{amkurth-Thesis/chapter3/figures/high_low_risk_full.pdf}
  \caption[Full Model, Mean Bootstrap Probability Estimates]{Full model: Mean bootstrap probability for each birth-weight category (C1–C11) in the high- and low-risk subgroups. Error bars show the 95 \% percentile interval over the $B$ bootstrap trees.}
  \label{fig:high-low-risk-full}
\end{figure}

% ─────────────────────────────────────────────────────────────────────────
% Full - NBW (C11 removed, 10 categories)
\begin{figure}[H]
  \centering
  \includegraphics[width=\textwidth]{amkurth-Thesis/chapter3/figures/high_low_risk_full_without_NBW.pdf}
  \caption[Full Model Without NBW Column, Mean Bootstrap Probability Estimates]{Full model (excluding the NBW column): Probability estimates and intervals are as in Fig.~\ref{fig:high-low-risk-full}.}
  \label{fig:high-low-risk-full-w/o-NBW}
\end{figure}

% ─────────────────────────────────────────────────────────────────────────
% LBW-only (10 categories)
\begin{figure}[H]
  \centering
  \includegraphics[width=\textwidth]{amkurth-Thesis/chapter3/figures/high_low_risk_small.pdf}
  \caption[LBW-Only, Mean Bootstrap Probability Estimates]{LBW-Only model (10 categories):  Mean bootstrap probability for each birth-weight category in the two risk groups with 95\% percentile intervals.}
  \label{fig:high-low-risk-lbw}
\end{figure}

% Depth distributions — Full model
\begin{figure}[H]
    \centering
    \includegraphics[width=1.0\textwidth]{amkurth-Thesis/chapter3/figures/boot/depth_distributions.png}
    \caption[Full Model, Distribution of Variable Depths]{Full model: Distribution of variable depths across the ensemble. Each panel shows a histogram indicating how frequently a given variable appears at each tree depth, where depth 0 corresponds to the root node. Variables closer to the root are generally more important in the model.}
    \label{fig:var-depth-distributions-full}
\end{figure}

% Depth distributions — LBW-only model
\begin{figure}[H]
    \centering
    \includegraphics[width=1.0\textwidth]{amkurth-Thesis/chapter3/figures/boot/depth_distributions_2.png}
        \caption[LBW-Only Model, Distribution of Variable Depths]{LBW-Only model: Distribution of variable depths across the ensemble. Each panel shows a histogram indicating how frequently a given variable appears at each tree depth, where depth 0 corresponds to the root node. Variables closer to the root are generally more important in the model.}
    \label{fig:var-depth-distributions-lbw}
\end{figure}


% TABLES -------------------------------------------------------
\begin{table}[H]
\centering
\caption[Comparison of Variable Importance in Full Model and LBW-Only Model]{Comparison of Variable Importance in Full Model and LBW-Only Model}
\label{tab:var_comparison}
\begin{tabular}{lcc}
\toprule
& \multicolumn{1}{c}{Full Model} & \multicolumn{1}{c}{LBW-Only Model} \\
\cmidrule(lr){2-2} \cmidrule(lr){3-3}
\multicolumn{3}{l}{\textit{Initial Split Variable}} \\
\midrule
mrace15 & 1.0000 & 1.0000 \\
\midrule
\multicolumn{3}{l}{\textit{Variable Frequency}} \\
\midrule
sex & 1.0000 & 1.0000 \\
dmar & 1.0000 & 0.0915 \\
mrace15 & 1.0000 & 1.0000 \\
mager & 1.0000 & 0.9741 \\
precare5 & 1.0000 & 0.6351 \\
cig\_0 & 1.0000 & 0.0196 \\
meduc & 0.3794 & 0.0000 \\
\bottomrule
\multicolumn{3}{p{0.8\textwidth}}{\textit{Note:} The table shows variable frequency (proportion of trees containing each variable) and initial split variable (normalized measure of predictive contribution) for both models.} \\
\end{tabular}
\end{table}

% ---------------------------  FULL MODEL ---------------------------
\begin{table}[H]
\centering
\caption[Mean Probability Estimates With 95\% Bootstrap Percentile Intervals for for High- and Low-Risk Birth-Weight Subgroups Under the Full Model]{Mean probability estimates $\hat{\pi}_{i,k}$ (with 95\% bootstrap percentile intervals) for high‑ and low‑risk birth‑weight subgroups under the \emph{full} model.}
\label{tab:pi_mean_full}
\begin{tabular}{ccccccc}
\toprule
\multirow{2}{*}{Category $k$} & \multicolumn{3}{c}{High‑risk} & \multicolumn{3}{c}{Low‑risk} \\
\cmidrule(lr){2-4}\cmidrule(lr){5-7}
 & $\hat{\pi}$ & 2.5\% & 97.5\% & $\hat{\pi}$ & 2.5\% & 97.5\% \\
\midrule
1  & 1.90 \% & 0.0180 & 0.0202 & 0.80 \% & 0.0075 & 0.0088 \\
2  & 1.72 \% & 0.0163 & 0.0181 & 0.86 \% & 0.0080 & 0.0096 \\
3  & 1.58 \% & 0.0150 & 0.0166 & 0.88 \% & 0.0080 & 0.0100 \\
4  & 1.63 \% & 0.0156 & 0.0170 & 0.90 \% & 0.0083 & 0.0102 \\
5  & 1.69 \% & 0.0162 & 0.0176 & 0.98 \% & 0.0090 & 0.0112 \\
6  & 1.61 \% & 0.0154 & 0.0168 & 0.91 \% & 0.0087 & 0.0098 \\
7  & 1.63 \% & 0.0157 & 0.0170 & 0.93 \% & 0.0088 & 0.0101 \\
8  & 1.78 \% & 0.0172 & 0.0185 & 1.08 \% & 0.0102 & 0.0117 \\
9  & 1.32 \% & 0.0126 & 0.0138 & 0.81 \% & 0.0075 & 0.0090 \\
10 & 1.52 \% & 0.0146 & 0.0158 & 0.96 \% & 0.0089 & 0.0106 \\
11 & 83.62 \% & 0.8330 & 0.8391 & 90.92 \% & 0.9026 & 0.9132 \\
\bottomrule
\multicolumn{7}{p{0.9\textwidth}}{\textit{Note:} $\hat{\pi}$ values are reported as percentages; percentile‑limit columns remain on the \([0,1]\) scale.  Estimates are based on $B$ bootstrap replicates.}\\
\end{tabular}
\end{table}

\begin{table}[H]
\centering
    \caption[Mean Probability Estimates With 95\% Bootstrap Percentile Intervals for for High- and Low-Risk Birth-Weight Subgroups Under the LBW-Only Model]{Mean probability estimates $\hat{\pi}_{i,k}$ (with 95\% bootstrap percentile intervals) for high‑ and low‑risk birth‑weight subgroups under the \emph{LBW‑only} model.}
\label{tab:pi_mean_lbw}
\begin{tabular}{ccccccc}
\toprule
\multirow{2}{*}{Category $k$} & \multicolumn{3}{c}{High‑risk} & \multicolumn{3}{c}{Low‑risk} \\
\cmidrule(lr){2-4}\cmidrule(lr){5-7}
 & $\hat{\pi}$ & 2.5\% & 97.5\% & $\hat{\pi}$ & 2.5\% & 97.5\% \\
\midrule
1  & 12.28 \% & 0.1143 & 0.1267 & 8.63 \% & 0.0802 & 0.0928 \\
2  & 10.61 \% & 0.1014 & 0.1090 & 9.73 \% & 0.0911 & 0.1030 \\
3  & 9.75 \% & 0.0935 & 0.1003 & 9.87 \% & 0.0930 & 0.1041 \\
4  & 9.90 \% & 0.0965 & 0.1018 & 10.18 \% & 0.0974 & 0.1064 \\
5  & 10.43 \% & 0.1014 & 0.1069 & 10.96 \% & 0.1054 & 0.1145 \\
6  & 9.74 \% & 0.0948 & 0.1011 & 10.16 \% & 0.0970 & 0.1050 \\
7  & 9.83 \% & 0.0953 & 0.1033 & 10.30 \% & 0.0973 & 0.1082 \\
8  & 10.57 \% & 0.1031 & 0.1088 & 11.42 \% & 0.1078 & 0.1218 \\
9  & 8.01 \% & 0.0776 & 0.0837 & 8.81 \% & 0.0835 & 0.0933 \\
10 & 8.87 \% & 0.0858 & 0.0941 & 9.94 \% & 0.0940 & 0.1053 \\
\bottomrule
\multicolumn{7}{p{0.9\textwidth}}{\textit{Note:} $\hat{\pi}$ values are reported as percentages; percentile‑limit columns remain on the \([0,1]\) scale.  Estimates are based on $B$ bootstrap replicates.}\\
\end{tabular}
\end{table}
